\chapter{研究基础与数据边界}

本报告将事实、估计和观点分层呈现。事实层来自官方业绩预告、Q1财务包、行情和行业数据；估计层来自经审计的情景模型与有限的原始研报/公开共识；观点层只用于定义验证顺序，不替代事实或外部锚。

\noindent
\begin{tabularx}{\textwidth}{L{2.1cm}L{4.0cm}X}
\toprule
\textbf{证据层级} & \textbf{用途} & \textbf{边界} \\
\midrule
一级：官方披露与市场数据 & 半年报预告、Q1财务、当前价格、行业指数和资金表。 & 反映已发生或可观察数据，不证明未来盈利持续性。 \\
二级：原始券商PDF与可审计共识快照 & 盈利预测、目标价、估值方法和可比框架。 & 仅当前、可比较字段可获得正权重。 \\
三级：摘要、转载、陈旧目标与用户文章 & 识别市场叙事、历史区间与待验证假设。 & 权重为零，不能被呈现为市场一致预期。 \\
本机构情景 & 估值区间、概率权重、行动门槛与失效条件。 & 明确标注为本机构假设，需由下一轮证据验证。 \\
\bottomrule
\end{tabularx}

全量候选的逐票公式、输入变量、方法理由、倍数理由、概率权重和证据路径均在附录B保留。这使主报告保持可读，同时不牺牲模型的可复算性。

\noindent
\begin{disclosurebox}[使用边界]
“主力资金”沿用公开数据源按成交单大小分类的统计口径，不代表对资金最终受益人、机构身份或一致行动关系的确认。行业阶段使用截至2026年7月14日的完整31行业日报和截至2026年7月10日的完整周报；连续流入公开表仅30/112，已按部分证据处理。全部目标价均为概率情景研究框架，不构成收益承诺或证券买卖建议。
\end{disclosurebox}
